 \chapter{三角形与四边形}\label{ch:g6:shapes}

三角形和四边形是几何学中最基本的图形。本章的目标是
认识它们的特殊类别 --- 等腰三角形、
等边三角形、直角三角形；长方形、菱形、正方形 --- 通过
它们的\emph{定义性质}，并
能用直尺和圆规精确地作出这些图形。

\section{三角形}

\begin{definition}[特殊三角形]\label{def:g6:shapes:triangles}
三角形 $ABC$ 有三个顶点和三条边。若满足以下条件之一，则称其为：
\begin{itemize}
  \item \emph{等腰三角形}\index{等腰三角形}：在 $A$ 处，当
        $AB = AC$（两条边相等）；
  \item \emph{等边三角形}\index{等边三角形}：当
        $AB = BC = CA$（三条边相等）；
  \item \emph{直角三角形}\index{直角三角形}：在 $A$ 处，当角
        $\widehat{BAC}$ 为直角。与直角相对的边 $[BC]$
        称为\emph{斜边}。
\end{itemize}
\end{definition}

\begin{omfigure}
\begin{tikzpicture}[scale=0.85]
  % isosceles
  \draw[thick] (0,0) -- (2.4,0) -- (1.2,2.1) -- cycle;
  \draw (0.55,0.96) -- (0.72,0.86);
  \draw (1.85,0.96) -- (1.68,0.86);
  \node[below, font=\small] at (1.2,-0.15) {等腰};
  \node[above, font=\small] at (1.2,2.1) {$A$};
  \node[below left, font=\small] at (0,0) {$B$};
  \node[below right, font=\small] at (2.4,0) {$C$};
  % equilateral
  \begin{scope}[xshift=4.6cm]
  \draw[thick] (0,0) -- (2.4,0) -- (1.2,2.08) -- cycle;
  \draw (1.13,0.05) -- (1.27,-0.05);
  \draw (0.55,0.99) -- (0.72,0.89);
  \draw (1.85,0.99) -- (1.68,0.89);
  \node[below, font=\small] at (1.2,-0.25) {等边};
  \end{scope}
  % right
  \begin{scope}[xshift=9.2cm]
  \draw[thick] (0,0) -- (2.6,0) -- (0,1.8) -- cycle;
  \draw[thin] (0.3,0) -- (0.3,0.3) -- (0,0.3);
  \node[below, font=\small] at (1.3,-0.15) {直角};
  \node[omThm, font=\small, rotate=-34] at (1.5,1.05) {斜边};
  \end{scope}
\end{tikzpicture}

{\small 三种特殊三角形。小短划表示相等的
边；方块表示直角。}
\end{omfigure}

\begin{method}[由三边作三角形]\label{met:g6:shapes:construct}
作 $ABC$，使 $AB = 5$ cm，$AC = 4$ cm，$BC = 3$ cm：
\begin{enumerate}
  \item 用直尺画线段 $[AB]$，其长为 $5$ cm；
  \item 以 $A$ 为圆心、$4$ cm 为半径画圆弧
        （即 $4$ cm 是圆上各点到 $A$ 的距离）；
  \item 以 $B$ 为圆心、$3$ cm 为半径画圆弧；
  \item 两段圆弧的交点即为点 $C$；将其与 $A$、
        $B$ 相连。
\end{enumerate}
\end{method}

\begin{omfigure}
\begin{tikzpicture}[scale=1.0]
  \coordinate (A) at (0,0);
  \coordinate (B) at (5,0);
  \coordinate (C) at (3.2,2.4);
  \draw[thick] (A) -- (B);
  \draw[thick, omDef] (A) -- (C);
  \draw[thick, omThm] (B) -- (C);
  \draw[omDef, thin] (C) arc[start angle=36.9, delta angle=22, radius=4];
  \draw[omDef, thin] (C) arc[start angle=36.9, delta angle=-22, radius=4];
  \draw[omThm, thin] (C) arc[start angle=126.9, delta angle=22, radius=3];
  \draw[omThm, thin] (C) arc[start angle=126.9, delta angle=-22, radius=3];
  \fill (A) circle (1.5pt) node[below left, font=\small] {$A$};
  \fill (B) circle (1.5pt) node[below right, font=\small] {$B$};
  \fill (C) circle (1.5pt) node[above, font=\small] {$C$};
  \node[omDef, font=\small] at (1.2,1.5) {$4$ cm};
  \node[omThm, font=\small] at (4.6,1.4) {$3$ cm};
  \node[below, font=\small] at (2.5,0) {$5$ cm};
\end{tikzpicture}

{\small 用圆规作三角形：顶点 $C$ 是
两段圆弧的交点。（在\cref{ch:g7:triangles}中
将学习何时能作出这样的三角形。）}
\end{omfigure}

\section{四边形}

\begin{definition}[特殊四边形]\label{def:g6:shapes:quadrilaterals}
四边形有四个顶点和四条边。若它是：
\begin{itemize}
  \item \emph{长方形}\index{长方形}：当其四个角
        均为直角；
  \item \emph{菱形}\index{菱形}：四条边
        长度均相等；
  \item \emph{正方形}\index{正方形}：同时满足两者，即四个角都是直角
        \emph{且}四条边相等。
\end{itemize}
\end{definition}

\begin{proposition}[对角线的性质]\label{prop:g6:shapes:diagonals}
在这些四边形中，两条对角线相交于它们的
公共中点，并且：
\begin{itemize}
  \item 在长方形中，对角线\emph{长度相等}；
  \item 在菱形中，对角线\emph{互相垂直}；
  \item 在正方形中，两者同时成立。
\end{itemize}
\end{proposition}
\admitted

\begin{omfigure}
\begin{tikzpicture}[scale=0.8]
  % rectangle, right-angle marks drawn inward at each corner
  \draw[thick] (0,0) rectangle (3,1.9);
  \draw[omThm] (0,0) -- (3,1.9);
  \draw[omThm] (0,1.9) -- (3,0);
  \foreach \x/\y/\dx/\dy in {0/0/1/1, 3/0/-1/1, 0/1.9/1/-1, 3/1.9/-1/-1}
    \draw[thin] (\x+0.22*\dx,\y) -- (\x+0.22*\dx,\y+0.22*\dy)
      -- (\x,\y+0.22*\dy);
  \node[below, font=\small] at (1.5,-0.2) {长方形};
  % rhombus
  \begin{scope}[xshift=5.1cm, yshift=0.95cm]
  \draw[thick] (0,0) -- (1.6,1.0) -- (3.2,0) -- (1.6,-1.0) -- cycle;
  \draw[omThm] (0,0) -- (3.2,0);
  \draw[omThm] (1.6,1.0) -- (1.6,-1.0);
  \draw[thin] (1.6,0) ++(0.18,0) -- ++(0,0.18) -- ++(-0.18,0);
  \node[below, font=\small] at (1.6,-1.2) {菱形};
  \end{scope}
  % square, with its right-angle marks too
  \begin{scope}[xshift=9.8cm]
  \draw[thick] (0,0) rectangle (1.9,1.9);
  \draw[omThm] (0,0) -- (1.9,1.9);
  \draw[omThm] (0,1.9) -- (1.9,0);
  \foreach \x/\y/\dx/\dy in {0/0/1/1, 1.9/0/-1/1, 0/1.9/1/-1,
                             1.9/1.9/-1/-1}
    \draw[thin] (\x+0.22*\dx,\y) -- (\x+0.22*\dx,\y+0.22*\dy)
      -- (\x,\y+0.22*\dy);
  \node[below, font=\small] at (0.95,-0.2) {正方形};
  \end{scope}
\end{tikzpicture}

{\small 三种特殊四边形及其对角线（红色）。
正方形兼具长方形与
菱形的性质。}
\end{omfigure}

\begin{remark}
正方形自动是长方形（它有四个直角），
也自动是菱形（四条边相等）。因此，“长方形”并不
意味着“不是正方形”！在数学中，较特殊的图形
属于较一般的那一类。
\end{remark}

\begin{example}[由性质识别图形]\label{ex:g6:shapes:recognize}
一个四边形的四条边长均为 $4$ cm，且有一个直角。
它是什么图形？四条边相等的四边形是菱形。在菱形中，对角
相等，四个角之和为 $360^\circ$；若其中一个角为直角，
则四个角都是直角（这一点将在
\cref{ch:g7:central}中借助平行四边形证明）。因此，该图形是正方形。
\end{example}

\begin{method}[作图程序]\label{met:g6:shapes:program}
一个\emph{作图程序}是一份指令清单，任何人都能
按照它准确复现一个图形。每条指令单独占一行，
创建每个点时都要为其命名，并且只使用以下指令：“画线段
/ 直线 / 圆\dots”，“将点放在\dots\ 相交处”。检验
你的程序：逐字照做。
\end{method}

\begin{example}\label{ex:g6:shapes:program}
正方形 $ABCD$ 的作图程序，边长为 $4$ cm：
\begin{enumerate}
  \item 画线段 $[AB]$，使 $AB = 4$ cm。
  \item 作 $(AB)$ 的垂线，过 $A$，并在其上取点
        $D$，使 $4$ cm 为其到 $A$ 的距离（在直线的一侧任选，并始终保持同一侧）。
  \item 作 $(AB)$ 的垂线，过 $B$，并在其上取点
        $C$，使 $4$ cm 为其到 $B$ 的距离，且与 $D$ 位于同一侧。
  \item 连接 $C$ 与 $D$。
\end{enumerate}
用直尺检验：$DC$ 的长度也应为 $4$ cm。
\end{example}

\section{练习}

\begin{exercise}[$\star$]\label{exo:g6:shapes:1}
作三角形 $ABC$，使 $AB = 6$ cm，$AC = 5$ cm，且
$BC = 4$ cm，按照\cref{met:g6:shapes:construct}作图。
\end{exercise}

\begin{exercise}[$\star$]\label{exo:g6:shapes:2}
作等腰三角形 $DEF$，使 $DE = DF = 5$ cm，且
$EF = 3$ cm。该三角形中哪两个角看起来相等？测量它们
以检验。
\end{exercise}

\begin{exercise}[$\star$]\label{exo:g6:shapes:3}
仅用直尺和圆规作边长为 $4$ cm 的
等边三角形，并测量其中一个角。
\end{exercise}

\begin{exercise}[$\star$]\label{exo:g6:shapes:4}
根据各三角形的数据对其分类（等腰、等边、
直角三角形 --- 可能有多个答案）：
(a) 三边长分别为 $7$、$7$、$7$； \
(b) 三边长分别为 $5$、$5$、$8$； \
(c) 三个角分别为 $90^\circ$、$45^\circ$、$45^\circ$； \
(d) 三个角分别为 $60^\circ$、$60^\circ$、$60^\circ$。
\end{exercise}

\begin{exercise}[$\star$]\label{exo:g6:shapes:5}
直角三角形中，哪条边是斜边？画一个三角形
$MNP$，使其在 $N$ 处为直角，并指出它的斜边。哪个顶点
永远不与斜边相接？
\end{exercise}

\begin{exercise}[$\star$]\label{exo:g6:shapes:6}
判断下列说法是否正确，并用一句话说明每个答案的理由。
\begin{enumerate}
  \item 每个正方形都是长方形。
  \item 每个长方形都是正方形。
  \item 每个正方形都是菱形。
  \item 有一个直角的菱形是正方形。
\end{enumerate}
\end{exercise}

\begin{exercise}[$\star$]\label{exo:g6:shapes:7}
作一个长方形 $ABCD$，使 $AB = 5$ cm，$BC = 3$ cm。量出
它的两条对角线，你有什么发现？
\end{exercise}

\begin{exercise}[$\star$]\label{exo:g6:shapes:8}
作一个菱形，使其对角线长分别为 $6$ cm 和 $4$ cm。（用
\cref{prop:g6:shapes:diagonals}：先画两条对角线 ---
它们互相垂直，并在各自的中点处相交 --- 然后连接四个
端点。）
\end{exercise}

\begin{exercise}[$\star\star$]\label{exo:g6:shapes:9}
编写一个等腰三角形的作图程序，底边长为
$4$ cm，两腰长均为 $6$ cm，写法仿照
\cref{ex:g6:shapes:program}。与同学互换程序，
看对方能否不看你的图形而照着作出？
\end{exercise}

\begin{exercise}[$\star\star$]\label{exo:g6:shapes:10}
画一个圆心为 $O$ 的圆，并画出它的一条直径 $[AB]$。在圆上任取一
点 $M$，画出三角形 $AMB$，量出
$\widehat{AMB}$ 的度数。移动 $M$ 后再量一次。你有什么猜想？
（这一出人意料的事实在\cref{ch:g8:cosine}中给出证明。）
\end{exercise}

\begin{exercise}[$\star\star$]\label{exo:g6:shapes:11}
一个四边形的两条对角线等长，且在
公共中点处互相平分。试利用\cref{prop:g6:shapes:diagonals}，判断它必须是哪一种图形：
长方形、菱形还是正方形？
（注意：题目只给出了一个性质。）
\end{exercise}

\begin{exercise}[$\star\star\star$]\label{exo:g6:shapes:12}
在 $4 \times 4$ 的点阵网格上，共能画出多少个不同的正方形，
顶点都是格点 --- 还要计入“倾斜”的正方形？
至少画出一个倾斜的正方形，并说明它的四条边为什么相等。
\end{exercise}

\section{问题：对角线能告诉我们什么}

\begin{problem}[{周末问题 --- 仅由对角线重建每一种特殊
四边形，并数出
任意多边形的对角线条数}]
\label{pb:g6:shapes:1}
\Cref{prop:g6:shapes:diagonals}由图形出发，说明它的对角线
有什么性质。本题则\emph{反过来}玩这个游戏：已知
只有两条相交线段 --- 即将成为对角线的线段 --- 连接
它们的四个端点，会出现什么图形？你将建立一份完整的
“对角线词典”，理解\emph{为什么}
\cref{exo:g6:shapes:8}中的菱形作法行得通，
遇见一个尚无名称的图形，最后数出
边数远多于四条的多边形的对角线条数。

\medskip
\noindent\textbf{第一部分 --- 对角线词典。}
对每个问题，按照描述画出两条线段，
依次连接四个端点，判断所得的四边形是什么。
\begin{enumerate}
  \item 两条线段各长 $6$ cm，互相垂直，并在
        它们的公共中点处相交。量出所得图形的四条边和
        四个角，它是什么图形？
  \item 两条线段分别长 $8$ cm 和 $5$ cm，互相垂直，
        并在它们的公共中点处相交。会出现什么图形
        （见\cref{exo:g6:shapes:8}）？
  \item 两条线段各长 $7$ cm，二者等长但\emph{不}
        互相垂直，并在它们的公共中点处相交。
        量出各边与各角，你猜想这个图形
        是什么？
  \item 两条线段分别长 $8$ cm 和 $5$ cm，不互相垂直，
        并在它们的公共中点处相交。所得图形不属于
        \cref{def:g6:shapes:quadrilaterals}中列出的三种图形之一 ---
        请描述你所看到的（哪些边看起来相等？哪些边平行？）。
        这个四边形的名称将在
        \cref{ch:g7:central}中给出。
  \item 把第 1--4 题的结果总结成一本含四个条目的词典：
        对角线\emph{相等 / 垂直 / 二者兼有 /
        二者皆无}（两条对角线总在中点处相交）
        $\rightarrow$ 图形的名称（或描述）。
\end{enumerate}

\medskip
\noindent\textbf{第二部分 --- 作图为什么行得通。}
\begin{enumerate}[resume]
  \item 在第 2 题中，两条对角线把图形分成四个直角三角形，
        它们都以交点为一个顶点。每个三角形的两条直角边长分别是多少？
        说明为什么四个三角形彼此全等，
        并由此推出图形的四条边为什么相等 ---
        这正是
        \cref{exo:g6:shapes:8}的秘诀。
  \item 推出一种仅用圆规的检验方法：由对角线作出任一图形后，
        如何只借助圆规（不用直尺测量），
        检验它的四条边是否
        都相等？
  \item 编写一个作图程序
        （见\cref{met:g6:shapes:program}），作一个正方形，它的
        \emph{对角线}长为 $6$ cm --- 而不是边长。
  \item 在第 8 题的正方形中，边长比 $3$ cm（即对角线的一半）更长还是更短？
        量一量便知。答案的一半可以
        已经被证明：
        \emph{证明}：从一个顶点经中心走到相邻顶点，
        要绕行 $3 + 3$ cm，而直线段总比绕行路线短，
        因此边长小于
        $6$ cm。（证明它也大于 $3$ cm，
        并精确算出它的值，需要用到
        \cref{ch:g8:pythagoras}。）
  \item 词典的一个新条目：画线段 $[AC]$，
        长 $8$ cm；在其上取点 $H$，使 $3$ cm 恰为其到端点 $A$ 的长度；
        过 $H$ 作与之垂直的线段 $[BD]$，
        长 $5$ cm，并以 $H$ 为它的\emph{中点}。连接 $ABCD$，
        得到的图形是\emph{筝形}\index{筝形}（一条对角线平分
        另一条，且两者成直角）。
        利用第 6 题中的角上三角形，找出筝形中
        哪些边与哪些边分别相等。
\end{enumerate}

\medskip
\noindent\textbf{第三部分 --- 数一数对角线。}
\emph{对角线}是多边形中连接两个
不相邻顶点的线段。
\begin{enumerate}[resume]
  \item 画一个四边形、一个五边形（五个顶点）和一个
        六边形（六个顶点），并画出它们所有的对角线。
        数一数，各有多少条？
  \item 不画图，直接计算六边形的对角线条数。从一个顶点出发，
        有几条对角线？（六个顶点中，有几个顶点与它
        \emph{不}能用对角线相连？）
        把这个数乘以六个顶点，每条对角线
        恰好被计入两次，为什么？由此核对你在第 11 题中算出的
        结果。（同样的重复计数也出现在
        \cref{pb:g6:wholes:1}的握手问题中。）
  \item 利用第 12 题的方法，计算一个有 $10$ 个顶点的多边形的
        对角线条数，再计算有 $20$ 个顶点的多边形的对角线条数。
  \item 一个多边形恰有 $44$ 条对角线，它有多少个顶点？
        （试用第 13 题的方法，对几个
        候选数逐一检验。）
  \item 用最小的多边形检验你的方法：三角形有几条对角线，
        第 12 题的计数方法给出什么结果？
        请用一句话总结
        握手问题与对角线
        之间有何共同之处。
\end{enumerate}
\end{problem}
